% ---------------------------------------------------------------------
%  appendix-b-package.tex -- using the hpb.sty package (API v1.0)
% ---------------------------------------------------------------------
\section{Using the \texttt{hpb.sty} Package}
\label{app:package}

The notation of \cref{sec:notation} is distributed as a single \LaTeX{}
package, \texttt{hpb.sty} (version~1.1, MIT license). Every HPB diagram in
this paper is drawn with it. The point is standardization: glyph sizes,
stroke widths, colors, fonts and the placement of tags, labels and terminals
are fixed by the package, so diagrams drawn by different people look the
same and follow the layout rules of \cref{sec:notation}. Use it for papers,
reports and specifications; for web pages and documentation that change
often, HPB Markup (\cref{sec:tooling,app:markup}) draws the same notation
from text. This appendix summarizes the package interface.

\paragraph{Installation and requirements.}
Copy \texttt{hpb.sty} next to the main \texttt{.tex} file (on Overleaf,
upload it to the project root) and load it with \texttt{\textbackslash
usepackage\{hpb\}}. The package requires only TikZ, \texttt{xcolor} and
\texttt{xparse}, which are part of every current \TeX{} distribution; it
loads the TikZ libraries \texttt{positioning}, \texttt{calc},
\texttt{shapes.geometric}, \texttt{arrows.meta},
\texttt{decorations.pathmorphing}, \texttt{backgrounds} and \texttt{fit}.
It uses no engine-specific features, needs no shell escape and no external
files, and is intended for pdf\LaTeX{} as used by Overleaf and arXiv.
Compile once with the engine of your submission before relying on it.

\paragraph{Canvas and coordinates.}
Every diagram is an \texttt{hpbgraph} environment, a TikZ picture in which
the $x$ unit is one ordinal column and the $y$ unit is one lane. All
macros therefore take plain numbers: columns (integers, or decimals for
fine adjustment) and lanes, with $0$ for the main line, positive lanes
above it and negative lanes below it (\cref{sec:notation}). A child line
created by a line fork is simply another $y$ value (e.g., $-3.5$). The
environment accepts three keys: \texttt{col} (column width, default
1.1\,cm), \texttt{lane} (lane height, default 1.0\,cm; tags are anchored at
a fixed distance above their node, and 1.0\,cm is the smallest value that
keeps a tag clear of a branch on lane $+1$) and \texttt{legend} (width of
the legend and origin boxes, default 2.3\,cm). This paper sets one canvas
for all diagrams with \texttt{\textbackslash hpbset\{col=0.75cm, lane=1cm,
legend=1.7cm\}}, because one document uses one canvas. Any other key is passed to the
underlying \texttt{tikzpicture}; \texttt{\textbackslash hpbset} sets
defaults for a whole document. Only spacing scales with these keys; glyph
sizes are fixed in points by the notation. Plain TikZ commands can be mixed
freely with the package macros in the same coordinates. Internally the
package draws on six layers (background for gate and bundle bands, guides,
branch paths and connectors, main lines and streams, nodes, and a front
layer for tags and labels), so that the order of commands in the source
does not affect the result.

\begin{figure}[tbp]
  \centering
  \begin{minipage}[c]{0.41\linewidth}
\begin{lstlisting}[style=hpblatex,basicstyle=\ttfamily\scriptsize,
  moretexcs={hpbgate,hpbtarget}]
\begin{hpbgraph}[col=0.7cm, legend=1.5cm]
  % lanes and main line
  \hpblane{1}{imp}{Improv. (I)}{0}{8.6}
  \hpblane{-1}{iss}{Issue (X)}{0}{8.6}
  \hpbmainline{0}{6}
  \hpbmainline*{6}{8.6}
  \hpborigin{0}{Product A}
  % baselines
  \hpbrelease{1}{v2.0}
  \hpbrelease{5}{v2.0.1}
  \hpbplanned{8}{v2.1}{Jun}
  % branches
  \hpbbranch[label={I-40 Export}]
    {imp}{1}{1}{2}{6}{target}{8}
  \hpbbranch[label={X-77 Login}]
    {iss}{1}{-1}{3}{4}{merge}{5}
  % markers
  \hpbnow{6}{1}{-1}
  \hpbgate[1.5]{8}{1}{-1}{gate v2.1}
\end{hpbgraph}
\end{lstlisting}
  \end{minipage}\hfill
  \begin{minipage}[c]{0.57\linewidth}
    \centering
    % natural size (about 8.3cm wide): no scaling, glyphs at normative size
    \begin{hpbgraph}[col=0.7cm, legend=1.5cm]
      % lanes and main line
      \hpblane{1}{imp}{Improv. (I)}{0}{8.6}
      \hpblane{-1}{iss}{Issue (X)}{0}{8.6}
      \hpbmainline{0}{6}
      \hpbmainline*{6}{8.6}
      \hpborigin{0}{Product A}
      % baselines
      \hpbrelease{1}{v2.0}
      \hpbrelease{5}{v2.0.1}
      \hpbplanned{8}{v2.1}{Jun}
      % branches
      \hpbbranch[label={I-40 Export}]
        {imp}{1}{1}{2}{6}{target}{8}
      \hpbbranch[label={X-77 Login}]
        {iss}{1}{-1}{3}{4}{merge}{5}
      % markers
      \hpbnow{6}{1}{-1}
      \hpbgate[1.5]{8}{1}{-1}{gate v2.1}
    \end{hpbgraph}
  \end{minipage}
  \caption{A complete \texttt{hpb.sty} example (illustrative data): the
    source (left) and its rendering (right). Columns are ordinal
    (one event per column) and lanes follow \cref{sec:notation}. The issue X-77,
    observed in \texttt{v2.0}, is merged into the patch baseline
    \texttt{v2.0.1}; the improvement I-40 is open at NOW and targeted at the
    planned baseline \texttt{v2.1}, whose gate it keeps closed until it is
    verified, re-targeted or discarded.}
  \label{fig:package-example}
\end{figure}

\paragraph{Interface.}
\Cref{tab:api} lists the public macros and \cref{tab:api-keys} the options
of \texttt{\textbackslash hpbbranch}, the central macro: one call draws a
complete branch, i.e., the fork from its baseline, the start node with the
type letter, the run along the lane, the terminal and, for merged and
targeted branches, the connection to the closing or target baseline. Types
are written \texttt{imp}, \texttt{iss}, \texttt{task}, \texttt{lab} (or
\texttt{I}, \texttt{X}, \texttt{T}, \texttt{L}); an unknown type, terminal
or step state stops compilation with an explicit error message rather than
producing a silently wrong diagram. The color names (\texttt{hpbInk},
\texttt{hpbMain}, \texttt{hpbImp}, \texttt{hpbIss}, \texttt{hpbTask},
\texttt{hpbLab}, \texttt{hpbHw}, \texttt{hpbDisc}, \texttt{hpbNow},
\texttt{hpbGuide}, \texttt{hpbLate}) are part of the interface, so that
non-HPB figures, such as charts and state diagrams, can use the same palette.

\providecommand{\HPBmac}[1]{\texttt{\textbackslash #1}}
\providecommand{\HPBarg}[1]{\ensuremath{\langle}\textit{#1}\ensuremath{\rangle}}

\begin{table}[tbp]
  \centering
  \small
  \caption{Options of \texttt{\textbackslash hpbbranch} (all optional).}
  \label{tab:api-keys}
  \renewcommand{\arraystretch}{1.0}
  \begin{tabularx}{\linewidth}{@{}>{\raggedright\arraybackslash}p{3.0cm}
      >{\raggedright\arraybackslash}p{2.0cm}
      >{\raggedright\arraybackslash}X@{}}
    \toprule
    Key & Default & Effect \\
    \midrule
    \texttt{line=}\HPBarg{y} & \texttt{0}
      & $y$ of the line the branch forks from (child line or stream). \\
    \texttt{mergeline=}\HPBarg{y} & = \texttt{line}
      & $y$ of the line the branch merges into. \\
    \texttt{forkoff=}\HPBarg{cols} & \texttt{0.5}
      & Offset of the vertical drop from the fork baseline (\cref{sec:notation}). \\
    \texttt{mergeoff=}\HPBarg{cols} & \texttt{0.5}
      & Offset of the vertical rise before the merge baseline. \\
    \texttt{label=}\HPBarg{text} & none
      & Branch label (reference and title) at the start node. \\
    \texttt{labelpos=}\HPBarg{pos} & \texttt{above}
      & Label position, as for \HPBmac{hpblabel}. \\
    \texttt{letter=}\HPBarg{text} & type letter
      & Overrides the letter of the start node. \\
    \texttt{nostart} & --
      & Omits the start node. \\
    \texttt{faded} & --
      & Draws at 35\,\% opacity (former state, e.g., an old target). \\
    \bottomrule
  \end{tabularx}
\end{table}

\begin{table}[p]
  \centering
  \footnotesize
  \caption{Public macros of \texttt{hpb.sty} v1.1. Columns and lanes are
    plain numbers in diagram units; optional arguments in square brackets.
    \HPBarg{type}: \texttt{imp}, \texttt{iss}, \texttt{task}, \texttt{lab}
    (legend boxes also accept \texttt{main}, \texttt{hw}, \texttt{disc}).}
  \label{tab:api}
  \renewcommand{\arraystretch}{1.0}
  \begin{tabularx}{\linewidth}{@{}>{\raggedright\arraybackslash}p{3.0cm}
      >{\raggedright\arraybackslash}p{5.6cm}
      >{\raggedright\arraybackslash}X@{}}
    \toprule
    Macro & Arguments & Draws \\
    \midrule
    \multicolumn{3}{@{}l}{\textit{Canvas}} \\
    \texttt{hpbgraph} (env.) & \texttt{[col=, lane=, legend=, }\HPBarg{tikz}\texttt{]}
      & The canvas: $x$ = column, $y$ = lane. \\
    \HPBmac{hpbset} & \texttt{\{}\HPBarg{keys}\texttt{\}}
      & Document-wide defaults for the canvas keys. \\
    \HPBmac{hpblane} & \HPBarg{lane}\HPBarg{type}\HPBarg{label}\HPBarg{from}\HPBarg{to}
      & Dotted lane guide and legend box at its left end. \\
    \midrule
    \multicolumn{3}{@{}l}{\textit{Lines and baselines}} \\
    \HPBmac{hpbmainline}[\texttt{*}] & \HPBarg{from}\HPBarg{to}[\HPBarg{y}]
      & Main line, 3\,pt; starred: dashed (plan). \\
    \HPBmac{hpborigin} & \HPBarg{col}\HPBarg{label}[\HPBarg{y}]
      & Creation dot (not a baseline) and product box. \\
    \HPBmac{hpbrelease} & [\HPBarg{caption}]\HPBarg{col}\HPBarg{name}[\HPBarg{y}]
      & Released baseline with tag and arrow; optional caption (e.g., applies-to set). \\
    \HPBmac{hpbplanned}[\texttt{*}] & [\HPBarg{caption}]\HPBarg{col}\HPBarg{name}\HPBarg{date}[\HPBarg{y}]
      & Planned baseline, dashed tag, target date; starred: overdue (``!'' badge, red). \\
    \HPBmac{hpbhardware} & [\HPBarg{caption}]\HPBarg{col}\HPBarg{name}[\HPBarg{y}]
      & Hardware diamond, indigo tag, caption (default \emph{new ref.}). \\
    \HPBmac{hpbprocess} & [\HPBarg{caption}]\HPBarg{col}\HPBarg{name}[\HPBarg{y}]
      & Process baseline (square). \\
    \HPBmac{hpbstream}[\texttt{*}] & [\HPBarg{offset}]\HPBarg{fork col}\HPBarg{y}\HPBarg{to col}[\HPBarg{from y}]
      & Release stream leaving a baseline; starred: planned. \\
    \HPBmac{hpbmilestone} & \HPBarg{col}\HPBarg{label}[\HPBarg{y}]
      & Life-cycle milestone tick. \\
    \HPBmac{hpbseparator} & \HPBarg{y}\HPBarg{from}\HPBarg{to}
      & Thin rule between lane bands. \\
    \HPBmac{hpbcollapsed} & \HPBarg{col}\HPBarg{count}[\HPBarg{y}]
      & Count badge of collapsed history. \\
    \HPBmac{hpbeos} & \HPBarg{col}[\HPBarg{y}]
      & End-of-support bar labeled \textsf{EoS}. \\
    \HPBmac{hpblinefork} & \HPBarg{col}\HPBarg{from y}\HPBarg{to y}
      & Lineage connector to a child line. \\
    \midrule
    \multicolumn{3}{@{}l}{\textit{Branches and events}} \\
    \HPBmac{hpbbranch} & [\HPBarg{keys}]\HPBarg{type}\HPBarg{fork col}\HPBarg{lane}%
      \HPBarg{start col}\HPBarg{end col}\HPBarg{end}\HPBarg{merge col}
      & Complete branch; \HPBarg{end} $\in$ \{\texttt{merge}, \texttt{target},
        \texttt{open}, \texttt{discard}, \texttt{none}\}; \HPBarg{merge col}
        is empty for \texttt{open}, \texttt{discard}, \texttt{none}. \\
    \HPBmac{hpbtarget} & [\HPBarg{keys}]\HPBarg{from col}\HPBarg{lane}\HPBarg{merge col}
      & Additional dashed target continuation (re-targeting). \\
    \HPBmac{hpbstart} & \HPBarg{type}\HPBarg{col}\HPBarg{lane}\HPBarg{letter}
      & Start node with letter (stand-alone). \\
    \HPBmac{hpbnode} & \HPBarg{type}\HPBarg{col}\HPBarg{lane}
      & Plain event node. \\
    \HPBmac{hpbterminal} & \HPBarg{kind}\HPBarg{type}\HPBarg{col}\HPBarg{lane}
      & Terminal glyph (stand-alone). \\
    \HPBmac{hpbstep} & \HPBarg{type}\HPBarg{col}\HPBarg{lane}\HPBarg{state}
      & Step: \texttt{done}, \texttt{todo}, \texttt{pass}, \texttt{fail}. \\
    \HPBmac{hpbpoststep} & \HPBarg{type}\HPBarg{from col}\HPBarg{col}\HPBarg{lane}\HPBarg{state}
      & Step after closure, on a dotted continuation. \\
    \HPBmac{hpbdepends} & [\HPBarg{tikz}]\HPBarg{x,y}\HPBarg{x,y}\HPBarg{label}
      & Dependency connector (dotted, open head). \\
    \HPBmac{hpbpropagates} & [\HPBarg{tikz}]\HPBarg{x,y}\HPBarg{x,y}\HPBarg{label}
      & Propagation connector (dashed, filled head). \\
    \HPBmac{hpblabel} & [\HPBarg{pos}]\HPBarg{col}\HPBarg{lane}\HPBarg{text}
      & Label with white halo (default position \texttt{above}). \\
    \midrule
    \multicolumn{3}{@{}l}{\textit{Markers and annotations}} \\
    \HPBmac{hpbnow} & \HPBarg{col}\HPBarg{lane top}\HPBarg{lane bottom}
      & NOW marker across the given lanes. \\
    \HPBmac{hpbgate} & [\HPBarg{right}]\HPBarg{col}\HPBarg{lane top}\HPBarg{lane bottom}\HPBarg{label}
      & Gate band around a planned baseline. \\
    \HPBmac{hpbbundle} & \HPBarg{x,y}\HPBarg{x,y}\HPBarg{label}
      & Launch-bundle frame across lines. \\
    \HPBmac{hpbfit} & \HPBarg{picture}
      & Scales a picture down to the line width only if wider. \\
    \HPBmac{hpbcallout} & [\HPBarg{tikz}]\HPBarg{x,y}\HPBarg{text}\HPBarg{x,y}
      & Annotation box with leader line. \\
    \HPBmac{hpbaxis} & \HPBarg{from}\HPBarg{to}\HPBarg{y}
      & Ordinal time axis. \\
    \HPBmac{hpbtick} & \HPBarg{col}\HPBarg{text}\HPBarg{y}
      & Calendar tick under the axis. \\
    \bottomrule
  \end{tabularx}

\end{table}

\paragraph{Conventions that the package does not enforce.}
Two layout rules depend on the data and are left to the author (or to a
layout algorithm such as \cref{alg:layout}; the HPB Markup renderer
computes columns and lanes for you). First, each event must be given its own column;
the geometry
of a branch requires
$\textit{fork col} + \textit{forkoff} < \textit{start col} \le
\textit{end col} \le \textit{merge col} - 1$.
Second, when several branches leave or enter the same baseline on the same
side, the branch on the farther lane must be given the smaller
\texttt{forkoff} or \texttt{mergeoff}, so that the corners nest
(\cref{sec:notation}). For a figure in a one-column article, the natural width of
the picture should not much exceed the text width, so that, after
\texttt{\textbackslash hpbfit\{\dots\}} (which scales down only when
needed), the smallest text (6.5\,pt) stays above 6\,pt; a
practical estimate of the natural width is
$\textit{legend} + 0.35\,\mathrm{cm} + (\textit{last col} + 0.8) \cdot
\textit{col}$.

\paragraph{Example.}
\Cref{fig:package-example} shows a complete diagram and its source: a
product line with a released baseline \texttt{v2.0}, a defect fixed by the
patch baseline \texttt{v2.0.1}, an improvement targeted at the planned
baseline \texttt{v2.1}, the NOW marker and the gate of \texttt{v2.1}. The
diagram needs twelve macro calls, none of which contains a coordinate in
centimeters.

